\section{Methodology}\label{methodology}

How Flavunato computes concentration uncertainties from SIMNRA ion beam
analysis (IBA) data.

The uncertainty formalism follows the internal technical note of Diallo
\citep{diallo2025iba}, whose Eqs. (D1)--(D6) are reproduced below.
Everything that is not taken from that note --- numerical conventions,
the estimator for the detected counts, additional cross-section values,
fit-quality thresholds and the optional AUC rescaling --- is a
\textbf{program choice} and is labelled as such throughout. Such choices
are implementation decisions of this software, not recommendations of
the cited note.

\subsection{Overview}\label{overview}

Flavunato reads a SIMNRA \texttt{.xnra} file together with the
corresponding element-spectrum export (a tab-separated \texttt{.txt}
file), and reports, for each element in each layer of the fitted model,
a concentration and its uncertainty.

\textbf{From the \texttt{.xnra} file} (XML, IDF schema with SIMNRA
extensions) the program takes the fitted layer model --- layer
thicknesses in \(10^{15}\ \text{atoms/cm}^2\) and elemental
concentrations as atomic fractions --- together with the experimental
and simulated spectra, the energy calibration
\(E = \text{offset} + \text{gain}\times\text{channel}\), and the beam
and geometry metadata. Isotope-resolved concentrations are retained
separately from element totals; element totals are formed by summing
isotopes of the same element after stripping the mass number.

\textbf{From the element-spectrum export} the program takes the
per-element and per-isotope simulated partial spectra, channel by
channel. These carry the information the \texttt{.xnra} does not: how
the total simulated spectrum decomposes into individual elemental
contributions.

\textbf{From the user} the program takes the regions of interest (ROIs,
as channel ranges with an associated element list), the
charge--solid-angle uncertainty \(\Delta_{Q\Omega}\), the cross-section
uncertainty \(\Delta_{\sigma_A}\), and the choice of adjustment method.

The program does \textbf{not} perform the spectrum fit; that is done in
SIMNRA \citep{mayer1999simnra}. Flavunato consumes a converged fit and
propagates uncertainties onto it.

\subsection{Areal density and detected
counts}\label{areal-density-and-detected-counts}

The number of counts detected for element \(A\) is \citep[Eq.
1]{diallo2025iba}

\[
A_A = \frac{Q \cdot \Omega \cdot \int_0^{x_A^0} f \cdot N_A^{A_m B_n}(x) \cdot \sigma_A(x)\, dx}{\cos \alpha} \tag{D1}
\]

where \(Q\) is the collected charge, \(\Omega\) the detector solid
angle, \(f\) a normalization factor, \(N_A^{A_m B_n}(x)\) the atomic
concentration profile of element \(A\), \(\sigma_A(x)\) the reaction
cross section, and \(\alpha\) the incidence angle.

Inverting this relation gives the areal atomic density \citep[Eq.
2]{diallo2025iba}

\[
N_A = N_A^{A_m B_n} \cdot x = \frac{\cos \alpha \cdot A_A}{Q \cdot \Omega \cdot \sigma_A} \tag{D2}
\]

\textbf{Scope note.} Flavunato does not evaluate Eqs. (D1)--(D2).
Concentrations and layer thicknesses are read directly from the
converged SIMNRA layer model. Eqs. (D1)--(D2) are reproduced here
because they establish \emph{why} \(A_A\), \(Q\Omega\) and \(\sigma_A\)
are the three quantities that carry uncertainty into \(C_A\): they are
the terms through which the measurement enters the concentration. The
\emph{Uncertainty budget} below propagates precisely these three terms.

\subsection{\texorpdfstring{Detected counts
\(A_A\)}{Detected counts A\_A}}\label{detected-counts-a_a}

\(A_A\) is the count of element \(A\) at its reference peak, integrated
over the ROI. Flavunato offers two estimators.

\subsubsection{Spectrum method}\label{spectrum-method}

All results reported in this work use this method.

\textbf{Program choice.} \(A_A\) is the integrated \emph{simulated
partial spectrum} of element \(A\) over the ROI:

\[
A_A = \operatorname{round}\!\left( \sum_{i \in \text{ROI}} y_i^{\mathrm{sim},A} \right)
\]

where \(y_i^{\mathrm{sim},A}\) is the simulated contribution of element
\(A\) in channel \(i\), taken from the element-spectrum export. The ROI
is \textbf{inclusive at both ends}: channels \(i\) with
\(i_\mathrm{start} \le i \le i_\mathrm{end}\).

No experimental normalisation is applied. The statistical uncertainty
derived from this quantity (see \emph{Counting statistics}) therefore
reflects the \emph{predicted} counting statistics of the fitted model,
not the measured ones. This is a deliberate consequence of the
estimator: it is reproducible from the fit alone and separates element
\(A\) from the overlapping contributions of other elements in the ROI,
but it will underestimate scatter if the fit itself is poor. This is
assessed separately (see \emph{Fit quality}).

In isotope mode the same integral is taken over the isotope-resolved
partial spectrum instead of the element total.

Two diagnostics are recorded alongside \(A_A\) but do not enter any
uncertainty: the element fraction
\(f_A = \sum_{\text{ROI}} y^{\mathrm{sim},A} / \sum_{\text{ROI}} y^{\mathrm{sim,total}}\),
and the experimental ROI sum \(\sum_{\text{ROI}} y^{\mathrm{exp}}\).

\subsubsection{Approximate fallback --- not
recommended}\label{approximate-fallback-not-recommended}

\textbf{Program choice.} If no element-spectrum export has been loaded,
the program falls back to \(A_A = \lfloor R \cdot C_A \rfloor\), where
\(R\) is the experimental ROI total and \(C_A\) the atomic fraction.
This assumes counts partition in proportion to atomic concentration,
which ignores the cross-section and stopping differences between
elements; for a heavy element in a light matrix it can be wrong by
orders of magnitude.

\textbf{Always load the metadata file.} Results obtained with the
fallback are not suitable for quantification. Every exported row records
which estimator was used in its \texttt{aa\_method} field; rows must
read \texttt{spectrum}.

\subsection{Uncertainty budget}\label{uncertainty-budget}

\subsubsection{Combination}\label{combination}

The uncertainty on the concentration \(C_A\) combines three independent
contributions \citep[Eq. 3]{diallo2025iba}:

\[
\Delta C_A = \sqrt{(\Delta_{Q\Omega})^2 + (\Delta_{\sigma_A})^2 + (\Delta_{A_A})^2} \tag{D3}
\]

All three terms are \textbf{relative} uncertainties (dimensionless
fractions). \(\Delta_{Q\Omega}\) and \(\Delta_{\sigma_A}\) are constant
across all elements and all layers of a given measurement; only
\(\Delta_{A_A}\) varies, through the counting statistics of each
element.

\textbf{Program convention.} Flavunato reports each term twice: as the
relative quantity above, and as an \emph{absolute} quantity scaled by
the concentration, \(\Delta C_A^{\text{abs}} = C_A \cdot \Delta C_A\).
The absolute form is the uncertainty on the concentration expressed in
the same units as the concentration itself (percentage points of atomic
fraction); the relative form is the one to quote when comparing
elements. Exported column names carry the suffix \texttt{\_rel\_percent}
or \texttt{\_percent} respectively.

\subsubsection{\texorpdfstring{Charge and solid angle,
\(\Delta_{Q\Omega}\)}{Charge and solid angle, \textbackslash Delta\_\{Q\textbackslash Omega\}}}\label{charge-and-solid-angle-delta_qomega}

\(\Delta_{Q\Omega}\) covers the combined uncertainty on the analysis
charge and the detector solid angle. Following \citep{diallo2025iba}, it
is estimated as the standard deviation of the set of ratios between the
simulated Particle\(\cdot\)sr values and the corresponding raw
(measured) Particle\(\cdot\)sr values obtained during normalisation. The
estimate is valid for samples originating from the same analysis
session, and is a property of the experimental setup rather than of any
element.

\textbf{Program choice.} \(\Delta_{Q\Omega}\) is a user input, expressed
in percent (default 3 \%). It is \emph{not} derived from the
\texttt{.xnra} file: the program does not compute the ratio statistics
above, which must be formed by the analyst across the session's samples
and entered manually. The spectrum-plotting module does read the beam
fluence and the detector solid angle from the file --- converting
millisteradians to steradians where the unit attribute requires it ---
and displays their product as Particle\(\cdot\)sr, which is the quantity
entering the ratios; but this value is used for display only and never
feeds the uncertainty calculation.

\subsubsection{\texorpdfstring{Counting statistics,
\(\Delta_{A_A}\)}{Counting statistics, \textbackslash Delta\_\{A\_A\}}}\label{counting-statistics-delta_a_a}

The statistical contribution is Poissonian \citep[Eq. 4]{diallo2025iba}:

\[
\Delta_{A_A} = \frac{\sqrt{A_A}}{A_A} = \frac{1}{\sqrt{A_A}} \tag{D4}
\]

\textbf{Program choice.} When \(A_A = 0\) the term is undefined and is
set to ``not a number''; the combination of Eq. (D3) then proceeds over
the two remaining terms only. See \emph{Assumptions and limitations} for
why this case requires attention.

\subsubsection{\texorpdfstring{Cross sections,
\(\Delta_{\sigma_A}\)}{Cross sections, \textbackslash Delta\_\{\textbackslash sigma\_A\}}}\label{cross-sections-delta_sigma_a}

\(\Delta_{\sigma_A}\) depends on the nuclear reaction, not on the
element as such. The note gives four values \citep{diallo2025iba}:

{\def\LTcaptype{none} % do not increment counter
\begin{longtable}[]{@{}
  >{\raggedright\arraybackslash}p{(\linewidth - 4\tabcolsep) * \real{0.3333}}
  >{\raggedright\arraybackslash}p{(\linewidth - 4\tabcolsep) * \real{0.3333}}
  >{\raggedright\arraybackslash}p{(\linewidth - 4\tabcolsep) * \real{0.3333}}@{}}
\toprule\noalign{}
\begin{minipage}[b]{\linewidth}\raggedright
Reaction
\end{minipage} & \begin{minipage}[b]{\linewidth}\raggedright
\(\Delta_{\sigma_A}\)
\end{minipage} & \begin{minipage}[b]{\linewidth}\raggedright
Source
\end{minipage} \\
\midrule\noalign{}
\endhead
\bottomrule\noalign{}
\endlastfoot
\(\mathrm{M}x(\alpha,\alpha)\mathrm{M}x\) & 0 \% & Diallo 2025; reading
it as Rutherford is this program's \\
\(^{16}\mathrm{O}(d,\alpha_0)^{14}\mathrm{N}\) & 2 \% & Diallo 2025 \\
\(^{1}\mathrm{H}(^{4}\mathrm{He},{}^{1}\mathrm{H})^{4}\mathrm{He}\) & 2
\% & Diallo 2025 \\
\(^{12}\mathrm{C}(d,p_0)^{13}\mathrm{C}\) & 5 \% & Diallo 2025 \\
\end{longtable}
}

\textbf{Program additions.} The following are \emph{not} in the cited
note. They were added to the program's default table and their
provenance is not recorded in the source; they must be justified
independently before being relied upon.

{\def\LTcaptype{none} % do not increment counter
\begin{longtable}[]{@{}
  >{\raggedright\arraybackslash}p{(\linewidth - 4\tabcolsep) * \real{0.3333}}
  >{\raggedright\arraybackslash}p{(\linewidth - 4\tabcolsep) * \real{0.3333}}
  >{\raggedright\arraybackslash}p{(\linewidth - 4\tabcolsep) * \real{0.3333}}@{}}
\toprule\noalign{}
\begin{minipage}[b]{\linewidth}\raggedright
Reaction
\end{minipage} & \begin{minipage}[b]{\linewidth}\raggedright
\(\Delta_{\sigma_A}\)
\end{minipage} & \begin{minipage}[b]{\linewidth}\raggedright
Status
\end{minipage} \\
\midrule\noalign{}
\endhead
\bottomrule\noalign{}
\endlastfoot
\(^{15}\mathrm{N}(p,\alpha\gamma)^{12}\mathrm{C}\) & 3 \% & program
addition, provenance unverified \\
\(^{2}\mathrm{H}(^{4}\mathrm{He},{}^{2}\mathrm{H})^{4}\mathrm{He}\) & 3
\% & program addition, provenance unverified \\
\end{longtable}
}

\textbf{Program addition --- isotope defaults.} For isotope-resolved
work the program carries a per-isotope default table, each entry mapped
to one of the reactions above:

{\def\LTcaptype{none} % do not increment counter
\begin{longtable}[]{@{}
  >{\raggedright\arraybackslash}p{(\linewidth - 4\tabcolsep) * \real{0.3333}}
  >{\raggedright\arraybackslash}p{(\linewidth - 4\tabcolsep) * \real{0.3333}}
  >{\raggedright\arraybackslash}p{(\linewidth - 4\tabcolsep) * \real{0.3333}}@{}}
\toprule\noalign{}
\begin{minipage}[b]{\linewidth}\raggedright
Species
\end{minipage} & \begin{minipage}[b]{\linewidth}\raggedright
Default \(\Delta_{\sigma_A}\)
\end{minipage} & \begin{minipage}[b]{\linewidth}\raggedright
Basis
\end{minipage} \\
\midrule\noalign{}
\endhead
\bottomrule\noalign{}
\endlastfoot
H, \(^{1}\)H & 2 \% & Diallo 2025 \\
\(^{2}\)H & 3 \% & program addition, unverified \\
C, \(^{12}\)C & 5 \% & Diallo 2025 \\
\(^{13}\)C & 5 \% & program addition, unverified \\
N, \(^{14}\)N, \(^{15}\)N & 3 \% & program addition, unverified \\
O, \(^{16}\)O & 2 \% & Diallo 2025 \\
\(^{17}\)O, \(^{18}\)O & 2 \% & program addition, unverified \\
other isotopes of H, C, N, O & element value & via element key \\
any other element & 0 \% (Rutherford) & Diallo 2025 for
\(\mathrm{M}x(\alpha,\alpha)\mathrm{M}x\); limiting it to Rutherford
(typically heavy) targets is this program's reading \\
\end{longtable}
}

\(^{13}\)C, \(^{17}\)O and \(^{18}\)O reuse the value of the main
isotope, although they are usually measured through different reactions;
these values have no basis in the cited note.

Lookup proceeds from a user-supplied override, to the exact isotope key,
to the element key obtained by stripping the mass number, and finally to
the 0 \% Rutherford fallback.

All stable isotopes of H, C, N and O --- \(^{1}\)H, \(^{2}\)H,
\(^{12}\)C, \(^{13}\)C, \(^{14}\)N, \(^{15}\)N, \(^{16}\)O, \(^{17}\)O,
\(^{18}\)O --- are listed explicitly, so the element-key step is reached
only by a non-stable entry such as \(^{3}\)H or \(^{14}\)C, or by a
mistyped name. Such an entry silently takes the element's value
\textbf{and its reaction label}: \(^{3}\)H is reported with
\(\sigma = 2\,\%\) and the reaction
\(^{1}\)H(\(^{4}\)He,\(^{1}\)H)\(^{4}\)He, \(^{14}\)C with
\(\sigma = 5\,\%\) and \(^{12}\)C(d,p\(_0\))\(^{13}\)C.

The 0 \% fallback is reached only when the base element is absent from
the table altogether. The cited note gives 0 \% for the reaction
\(\mathrm{M}x(\alpha,\alpha)\mathrm{M}x\); treating absence from the
table as equivalent to ``measured by Rutherford scattering'' is this
program's reading, not a statement of the note. The fallback therefore
applies to any unlisted species regardless of how it was actually
measured: \(^{7}\)Li, for example, silently receives a 0 \%
cross-section uncertainty even when its cross section is not Rutherford.

\textbf{The global cross-section input is not consulted in isotope
mode.} For an ROI entry carrying a mass number, the value comes only
from the table above or from an explicit per-isotope override; the
single cross-section field that applies to element-mode entries is
ignored. This is a known limitation of the current version.

\subsection{Fit quality}\label{fit-quality}

\subsubsection{Reduced chi-squared}\label{reduced-chi-squared}

Agreement between simulated and experimental spectra is quantified by
\citep[Eq. 5]{diallo2025iba}

\[
\chi^2_\nu = \frac{1}{\nu} \sum_{i=1}^{N} \frac{\left(y_i^{\exp} - y_i^{\mathrm{sim}}\right)^2}{\sigma_i^2} \tag{D5}
\]

with \(\nu = N - p\) degrees of freedom, \(N\) the number of channels in
the ROI and \(p\) the number of fitted parameters.

\textbf{Program choices.} Two quantities the note leaves open are fixed
as follows.

\begin{itemize}
\tightlist
\item
  \(p = 1\), so \(\nu = N - 1\).
\item
  \(\sigma_i = \sqrt{\max\left(y_i^{\exp},\, 4\right)}\). The Poisson
  uncertainty \(\sqrt{y^{\exp}}\) is floored at 4 counts to avoid
  division by zero in empty channels. The floor suppresses the
  contribution of near-empty channels rather than letting them dominate;
  it has no effect in channels with at least 4 counts.
\end{itemize}

\subsubsection{Interpretation}\label{interpretation}

\textbf{Program choice.} The following thresholds are used to label fit
quality in the interface and in exported results. They are a reporting
convenience of the program and carry no statistical authority.

{\def\LTcaptype{none} % do not increment counter
\begin{longtable}[]{@{}ll@{}}
\toprule\noalign{}
Range & Label \\
\midrule\noalign{}
\endhead
\bottomrule\noalign{}
\endlastfoot
\(0.90 \le \chi^2_\nu \le 1.10\) & Perfect \\
\(1.10 < \chi^2_\nu \le 3.0\) & Great \\
\(3.0 < \chi^2_\nu \le 10.0\) & Not OK \\
\(\chi^2_\nu > 10.0\) & Very bad \\
\(\chi^2_\nu < 0.90\) & Too low \\
\end{longtable}
}

A value below 0.90 is flagged because it generally indicates
overestimated uncertainties or an over-parameterised fit, not a better
result.

\subsubsection{AUC ratio}\label{auc-ratio}

\textbf{Program addition.} Alongside \(\chi^2_\nu\) the program computes
a normalisation check, the ratio of integrated areas over the ROI:

\[
\mathrm{AUC} = \frac{\sum_{i \in \text{ROI}} y_i^{\mathrm{sim}}}{\sum_{i \in \text{ROI}} y_i^{\exp}}
\]

A value near unity indicates that simulated and measured yields are
correctly normalised, independently of whether the \emph{shape} agrees.
This quantity does not appear in the cited note.

\subsection{Adjusted uncertainty}\label{adjusted-uncertainty}

When \(\chi^2_\nu > 1\), the fit residuals exceed what counting
statistics alone predict, which may indicate underestimated
uncertainties or unmodelled systematic effects. In that case the
statistical term may be rescaled by \(\sqrt{\chi^2_\nu}\), giving
\citep[Eq. 6]{diallo2025iba}

\[
\Delta C_A^{\mathrm{adj}} = \sqrt{(\Delta_{Q\Omega})^2 + (\Delta_{\sigma_A})^2 + \left(\sqrt{\chi^2_\nu}\, \Delta_{A_A}\right)^2} \tag{D6}
\]

The note is explicit that the reduced chi-squared is \emph{not} an
independent contribution to the uncertainty budget --- it is a
goodness-of-fit indicator, not a physical source of uncertainty --- and
that the correction applies only when \(\chi^2_\nu\) significantly
exceeds unity. Otherwise the standard expression of Eq. (D3) is used.

\textbf{Program implementation.} The rescaling factor is computed as
\(\sqrt{\max\left(\chi^2_\nu,\, 1\right)}\), clamped at unity so that a
fit with \(\chi^2_\nu < 1\) can never \emph{reduce} the reported
uncertainty. The adjustment is applied only when the analyst selects it;
\textbf{the default is no adjustment}, so that quoted uncertainties
follow Eq. (D3) unless stated otherwise. The program also computes and
exports the adjusted values unconditionally, in columns suffixed
\texttt{\_chi2\_adjusted\_percent}, so that both can be compared without
recomputation.

The program does not decide what counts as ``significantly exceeds
unity''; that judgement remains with the analyst, informed by the
thresholds in \emph{Interpretation}.

\textbf{Program addition --- AUC adjustment.} A third option rescales
the statistical term by
\(\max\left(\mathrm{AUC},\, 1/\mathrm{AUC}\right)\) instead. This
penalises normalisation mismatch in either direction. It is an addition
of the program and has no basis in the cited note; it is reported for
comparison only.

\subsection{Depth profiles}\label{depth-profiles}

A depth profile tabulates, layer by layer, the concentration of each
element and its uncertainty.

\textbf{Program choice.} When an element has been quantified in more
than one ROI within the same layer, the program reports the row with the
\textbf{lowest} \(\chi^2_\nu\), on the reasoning that the best-fitting
region gives the most reliable determination. Each profile row carries
the layer number, the layer thickness in \(10^{15}\ \text{atoms/cm}^2\),
and per element the concentration (in atomic percent, 2 decimals) with
its uncertainty (3 decimals) taken from the column matching the selected
adjustment method.

This selection is applied regardless of whether a \(\chi^2\) adjustment
was chosen for the uncertainty itself; \(\chi^2_\nu\) therefore
influences the reported profile even under the default setting.

\subsection{Assumptions and
limitations}\label{assumptions-and-limitations}

The following constrain the interpretation of the results and should be
stated alongside any quoted uncertainty.

\textbf{Single \(\Delta_{Q\Omega}\) and \(\Delta_{\sigma_A}\) per run.}
The interface accepts one value of each, applied to every element and
layer. This matches the formalism, in which both are constant for a
measurement, but means that a sample requiring different cross-section
uncertainties per element (for instance RBS for a heavy element and NRA
for a light one in the same spectrum) must be processed in more than one
run, or through the isotope-resolved path.

\textbf{The statistical term is an ROI quantity, not a per-layer one.}
\(A_A\) is integrated over the ROI across the whole depth sampled by
that ROI. The same \(A_A\), and hence the same \(\Delta_{A_A}\), is
therefore applied to every layer in which the element appears. The
layer-to-layer variation in the reported absolute uncertainty comes only
from the concentration scaling described in \emph{Combination}, not from
layer-resolved counting statistics. Depth-resolved counting
uncertainties are outside the scope of this formalism.

\textbf{Background and peak overlap are not included.} Eq. (D3) combines
charge, cross-section and counting terms only. No variance is attributed
to background subtraction, to overlap between the partial spectra of
different elements within an ROI, or to the layer model itself (number
of layers, thickness assignment). Where ROIs overlap --- a common
situation for adjacent heavy elements --- the resulting correlation is
not propagated. Quoted uncertainties are therefore lower bounds on the
total uncertainty.

\textbf{The ROI end channel is treated inconsistently.} \(A_A\) includes
the final channel of the ROI, whereas the \(\chi^2_\nu\) summation
excludes it. The two quantities in a given result row thus describe
windows differing by one channel. The effect on \(\chi^2_\nu\) is of
order 0.1--1 \% for the ROI widths used here. The behaviour is retained
deliberately, since the results reported in this work were produced with
it; it is recorded as a known issue and is a candidate for correction in
a future version.

\textbf{A missing element name yields \(A_A = 0\) without warning.} If
an element named in an ROI is absent from the element-spectrum export
--- for example through a naming difference between the layer model and
the export --- the lookup returns zero counts silently. The statistical
term then becomes undefined and is dropped from Eq. (D3), so the
reported uncertainty collapses to
\(\sqrt{(\Delta_{Q\Omega})^2 + (\Delta_{\sigma_A})^2}\). Because this is
\emph{smaller} and visually cleaner than a correct value, the failure is
easy to miss. The diagnostic is the exported \texttt{element\_counts}
column, which reads 0 for any affected row; it should be checked before
results are used.

\textbf{\(\chi^2_\nu\) is restricted to the simulated channel range.} An
ROI extending beyond the last simulated channel is silently truncated to
that range, and one lying entirely beyond it returns an undefined
\(\chi^2_\nu\). In practice no useful ROI lies there, since the
simulation carries no signal beyond its own range.

\subsection{Validation}\label{validation}

Two independent checks are maintained.

\textbf{Agreement with the manual calculation.} The program's
implementation of Eq. (D3) is compared against an independent manual
spreadsheet calculation, following the method of \citep{diallo2025iba}.
Supplied with the same detected counts and the same \(\Delta_{Q\Omega}\)
and \(\Delta_{\sigma_A}\) per element, the program reproduces the manual
relative and absolute uncertainties for every layer and element to
within \(10^{-9}\); the largest observed deviation is at the level of
floating-point rounding (\(\sim 10^{-18}\)). This validates the
uncertainty formula and its combination, independently of how \(A_A\) is
obtained.

The counts themselves are \emph{not} expected to agree: the manual
calculation used measured counts, whereas the program integrates the
simulated partial spectrum (see \emph{Spectrum method}). The two are
different estimators of \(A_A\), so that comparison is reported as a
diagnostic and never as a pass/fail criterion.

\textbf{Regression against the thesis version.} A regression test
compares the current program output with the tagged thesis version
(\texttt{v1.0-thesis}) at an absolute tolerance of \(10^{-12}\) on every
numeric field; all later code changes reproduce it exactly.
